\documentclass[11pt]{article}

\usepackage[a4paper,margin=1in]{geometry}
\usepackage[T1]{fontenc}
\usepackage[utf8]{inputenc}
\usepackage{lmodern}
\usepackage{amsmath,amssymb}
\usepackage{physics}
\usepackage{hyperref}
\usepackage{enumitem}
\usepackage{setspace}
\usepackage{tcolorbox}
\tcbuselibrary{skins,breakable}

\newtcolorbox{gapbox}[1][]{enhanced,breakable,
  colback=orange!6!white,colframe=orange!60!black,
  fonttitle=\bfseries,title={#1},boxrule=0.6pt,arc=4pt}

\hypersetup{colorlinks=true,linkcolor=blue,citecolor=blue,urlcolor=blue}
\setstretch{1.1}

\title{Saturated Superfluid Vacuum IV\\
Gravity as Update-Capacity Gradient}
\author{Stig Norland}
\date{Draft}

\begin{document}
\maketitle

\begin{abstract}
This paper develops gravity in the Saturated Superfluid Vacuum (SSV)
framework as the bending of motion through gradients of local time delay.
It builds on Paper~II, where gravity enters the force sector as the
interference of the pressure waves that every massive breather radiates into
the vacuum, and on Paper~III, where physical time is identified with the
local rate of change of the order parameter $\Psi$ --- the medium's local
\emph{update capacity}.  The chain is short: the interference of radiated
waves consumes update capacity and so slows local time; that slowing is
uneven and therefore has a gradient; and any motion through the gradient is
bent toward the region of greatest delay.  That bending is gravity.  From a
single interference potential $\Phi$ follow free fall, gravitational time
dilation, light deflection, and redshift, each at leading order in
$\Phi/c^2$.  The treatment is weak-field throughout; recovery of the full
curved-spacetime description is left to Paper~VII-b.  Two sections --- the
formal account of update capacity and the quantitative refraction geometry
--- are marked as under development.
\end{abstract}

\newpage
\tableofcontents
\newpage

\section{Introduction}
\label{sec:intro}

Gravity is usually approached either as a field to be quantised or as the
curvature of a spacetime taken as fundamental.  This paper takes neither
route.  In the Saturated Superfluid Vacuum (SSV) the vacuum is a real
medium --- a saturated compressible superfluid described by an order
parameter $\Psi$ --- and the question is not how to quantise a pre-existing
gravitational field, but how gravitational behaviour arises at all if the
medium's local rate of change is finite and depends on how heavily the
medium is loaded.

Two earlier papers supply the foundation.  Paper~II places gravity in the
force sector: every massive defect is a breather, a localised oscillation of
$\Psi$, and it radiates a longitudinal pressure wave into the medium; the
time-averaged interference of two such waves is the gravitational
interaction, and matching it to the observed inverse-square law fixes
Newton's constant $G$ from the medium parameters.  Paper~III identifies
physical time with the local rate of change of $\Psi$: time is not a
backdrop but a measure of how fast the medium is updating.  The present
paper takes the mechanism introduced in Paper~II and develops it in its own
right, on the footing Paper~III provides.

The thesis is a single chain.  The interference of radiated waves slows the
local rate of change of $\Psi$ --- it imposes a time delay.  The delay is
uneven, strongest near a mass and fading with distance, so across space it
is a gradient.  Any motion through that gradient is bent toward the side
where time runs slower.  That bending is gravity.  A single potential
$\Phi$ summarises the delay, and from it follow free fall, gravitational
time dilation, light deflection, and redshift.

The treatment here is weak-field, exact at leading order in $\Phi/c^2$.  The
recovery of the full curved-spacetime description --- the effective metric,
the post-Newtonian coefficients, the dissolution of the quantum-gravity
problem --- is the subject of Paper~VII-b, to which \S\ref{sec:geometry}
hands off.  Galactic-scale applications of the same mechanism are treated in
Papers~VI-a and~VI-b.  Section~\ref{sec:capacity} formalises update
capacity; \S\ref{sec:picture} sets out the physical picture;
\S\S\ref{sec:interference}--\ref{sec:consequences} make it quantitative;
\S\ref{sec:refraction} develops the refraction geometry of path bending;
and \S\S\ref{sec:universal}--\ref{sec:remembers} treat universality, the
absence of a graviton, the bridge to emergent geometry, and the connection
to the recording medium of Paper~III.  Sections~\ref{sec:capacity}
and~\ref{sec:refraction} are still under development.

\section{Update Capacity and Local Change-Rate}
\label{sec:capacity}

\emph{[Section under development --- to be written.]}  This section will
formalise the notion of \emph{update capacity}: the vacuum medium has a
finite local rate of change, and internal interaction and translational
propagation draw on the same budget.  Proper time is identified with local
change-rate (Paper~III); where the medium is loaded with interaction, less
capacity remains for free propagation, local time runs slow, and the
resulting spatial gradient in capacity is what the rest of the paper
develops as gravity.

\section{The Physical Picture}
\label{sec:picture}

Gravity is the gradient of a time delay --- and that delay is built
by the interfering pressure waves that every mass radiates into the vacuum.
Three steps connect those two statements: the interference of waves slows
local time; that slowing is uneven, so across space it is a gradient; and
any motion through that gradient is bent toward the slower side.  The
bending is gravity.  This section sets out the chain;
\S\S\ref{sec:interference}--\ref{sec:refraction} make each step
quantitative.

In Paper~I every massive defect is a breather: a localised region of $\Psi$
oscillating at its Compton frequency.  A breather does not sit inertly in the
medium.  It radiates a longitudinal pressure wave --- a travelling density
modulation $\delta\rho(\mathbf{x},t)$ --- outward at the medium's sound speed
$c$, with amplitude falling as $1/r$ from spherical dilution.  The region
surrounding a mass is therefore not a patch of medium with some altered
static property; it is a region permeated by that radiated wave.

\emph{Interference slows local time.}
In SSV, time is not a backdrop through which things move; it is the local
rate at which $\Psi$ changes (Paper~III; \S\ref{sec:capacity}).  Any process
--- a clock, a particle's internal oscillation, a light wave --- advances by
changing $\Psi$, and in undisturbed vacuum it advances at the medium's full
rate.  A process embedded in a breather's radiated wave is not alone: the
medium around it is already oscillating with that wave, and the process must
run \emph{through} it.  The two interfere.  Through the medium's logarithmic
nonlinearity this interference has a definite, one-signed effect
(\S\ref{sec:interference}): it lowers the embedded process's rate of
advance, so that fewer of its cycles complete per unit coordinate time.
Because physical time \emph{is} that rate of advance, the process runs slow
--- and the more intense the radiated wave, the greater the slowing.  This
is the time delay.

\emph{The delay is uneven, so it has a gradient.}
The radiated wave dilutes as $1/r$, so the interference is most intense at
the mass and fades with distance.  Local time therefore runs slowest at the
mass and recovers the full vacuum rate far away.  The delay is not a uniform
offset added everywhere alike; it varies smoothly from place to place, and a
quantity that varies across space has a gradient --- here one that points
toward the mass and steepens as the mass is approached.

\emph{The gradient bends motion, and that is gravity.}
A gradient of time delay deflects everything that crosses it.  Anything that
moves is a wave or an oscillation advancing through the medium --- a light
ray carrying its front forward, a body carrying its Compton oscillation.
Where local time is uneven, the side of that advance nearer the mass keeps
slower pace than the side farther away; the front pivots, wheeling toward
the slow side, exactly as a wavefront refracts toward a slower medium or a
rank of marchers wheels toward the flank taking shorter steps.  The
trajectory curves toward the region of greatest time delay --- toward the
mass.  A light ray grazing the Sun bends this way: not by passing through
``denser matter'' but by advancing through the Sun's radiated wave, which
runs its near side slow.  Even a body at rest is internally a pulsating
breather, and its oscillation runs slow on the side facing the mass; that
tilt of its internal phase is itself momentum directed toward the slow side,
and the body falls.  An object never feels the distant mass --- it responds
only to the local slope of the time delay where it sits.  That slope, not
the mass, is the source of gravity; the mass merely shapes it.

The time-averaged interference is summarised by a single potential
$\Phi(\mathbf{x})$, and from it follow the four observed gravitational
phenomena, derived in turn below:
\begin{itemize}
  \item \textbf{free fall} --- a test mass drifts down the potential
        gradient (\S\ref{sec:refraction});
  \item \textbf{time dilation} --- a clock deeper in the potential ticks
        slow (\S\ref{sec:consequences});
  \item \textbf{light deflection} --- a light wave's front pivots in the
        potential gradient (\S\ref{sec:consequences},~\ref{sec:refraction});
  \item \textbf{gravitational redshift} --- a wave climbing out of the
        potential loses frequency (\S\ref{sec:consequences}).
\end{itemize}
The acceleration of free fall is $-\nabla\Phi$ and the fractional slowing of
clocks is $\Phi/c^2$.  The interference cross-term of
\S\ref{sec:interference} builds $\Phi$ from the radiated waves; the sections
after it read off the consequences.

\section{The Interference Cross-Term}
\label{sec:interference}

The mechanism is already contained in the LogSE.  The medium's nonlinear
self-interaction is the logarithmic term $b\ln(\rho/\rho_0)$ of the Paper~I
Lagrangian.  Consider a point of the medium reached by the radiated waves
of several breathers.  The local density is the saturation value
modulated by each wave,
\begin{equation}
    \frac{\rho(\mathbf{x},t)}{\rho_0} = 1 + \sum_i x_i(\mathbf{x},t),
    \qquad x_i \equiv \frac{\delta\rho_i}{\rho_0},
\end{equation}
and the logarithmic energy the medium carries there is governed by
$\ln(\rho/\rho_0)$.  Expanding,
\begin{equation}
    \ln\!\Big(1 + \sum_i x_i\Big)
      = \sum_i x_i
      - \frac{1}{2}\Big(\sum_i x_i\Big)^{\!2}
      + \cdots
      = \sum_i x_i
      - \frac{1}{2}\sum_i x_i^2
      - \sum_{i<j} x_i x_j
      + \cdots
    \label{eq:log_expansion}
\end{equation}
The single-wave terms $x_i$ and $x_i^2$ are each breather's own
self-energy.  The interaction between two distinct breathers $A$ and $B$ is
carried entirely by the cross term
\begin{equation}
    u_{AB} \;=\; -\,b\,\rho_0\,x_A\,x_B .
    \label{eq:cross_term}
\end{equation}
This is an \emph{interference} term: the product of two waves, not a
property of either alone.  Three features of \eqref{eq:cross_term} fix the
character of gravity.

\paragraph{It survives averaging only for coherent waves.}
Each $x_i$ oscillates at its breather's Compton frequency.  The observable
interaction is the time average $\langle u_{AB}\rangle = -b\rho_0\langle x_A
x_B\rangle$, and $\langle x_A x_B\rangle$ vanishes unless the two waves are
phase-correlated.  Incoherent waves do not gravitate toward one another;
coherent ones do.  All matter breathers are eigenmodes of the one medium and
lock in phase through it (a resonance condition; see Paper~II), so
$\langle x_A x_B\rangle > 0$ generically.

\paragraph{Its sign is fixed by the concavity of the logarithm.}
The cross term carries the minus sign of the $-\tfrac12(\sum x)^2$
expansion --- a direct consequence of $\ln$ being concave.  For
phase-correlated waves $\langle x_A x_B\rangle > 0$, so $\langle
u_{AB}\rangle < 0$: the interaction energy is lowered where the waves
overlap in phase.  \emph{Gravity is attractive because the medium's
nonlinearity is logarithmic, and the logarithm is concave.}

\paragraph{It is bilinear, so the potential goes as $1/r$.}
Equation~\eqref{eq:cross_term} is linear in each wave.  Breather $B$, in
its own localised volume, couples to breather $A$'s wave evaluated at $B$'s
location, where spherical dilution has reduced it to $x_A \propto m_A/r$;
its own wave has $x_B \propto m_B$.  The time-averaged interaction energy
is therefore
\begin{equation}
    \langle u_{AB}\rangle \;\propto\; -\,\frac{m_A m_B}{r},
\end{equation}
the Newtonian form.  The gravitational potential $\Phi$ is this interference
energy per unit test mass; the force is $-\nabla\Phi \propto m_A m_B/r^2$.
The inverse square is the gradient of a diluted wave, not a separately
postulated law.

The expansion fixes the \emph{structure} of $\Phi$ --- interference only,
attractive, $1/r$.  Its coefficient is fixed by the breather's acoustic
output; matching the result to the observed inverse-square law yields
Newton's constant, computed from the medium in Paper~II.

The same time-averaged interference of pulsating waves was first analysed
classically by Bjerknes~\cite{Bjerknes1906}, who obtained it as a
time-averaged force $\langle F\rangle = -\rho_0\langle\dot V_1\dot
V_2\rangle/(4\pi r^2)$ between two bodies pulsating in a compressible
medium.  Paper~II gives the historical credit, shows that $\langle\dot V_1
\dot V_2\rangle$ is exactly the cross-correlation $\langle x_A x_B\rangle$
above written in volume-velocity variables, and carries the result through
to $G$.  The derivation here is LogSE-native and needs no external
pulsator: the breathers are self-sustaining eigenmodes of $\Psi$.

\section{Time Dilation, Deflection, and Redshift}
\label{sec:consequences}

Paper~II fixes the interference potential of a source mass $M$,
\begin{equation}
    \Phi(r) = -\,\frac{GM}{r},
    \label{eq:Phi_newton}
\end{equation}
with $G$ computed from the medium; its gradient $-\nabla\Phi$ is free fall.
The same $\Phi$ --- with no further parameter --- gives the three remaining
gravitational phenomena.  Each is the interference cross-term of
\S\ref{sec:interference} read for a different probe.

\paragraph{Time dilation.}
A clock is an oscillator: any periodic physical process carries a rest
energy $mc^2$ and runs at $\omega = mc^2/\hbar$.  Placed in the wave
of a source mass, the clock's oscillation acquires the interference
cross-term with that wave; its energy is shifted by the interaction energy
$m\Phi$.  The fractional change in tick rate is therefore
\begin{equation}
    \frac{d\tau}{dt} = 1 + \frac{m\Phi}{mc^2} = 1 + \frac{\Phi}{c^2},
    \label{eq:time_dilation}
\end{equation}
the standard weak-field result; with $\Phi < 0$ in a well, the clock runs
slow.  This is not an independent derivation --- it is
\eqref{eq:Phi_newton} divided by $c^2$.  The slow clock is the time delay
of \S\ref{sec:picture}: each oscillation completes less physical change per
unit coordinate time where the interference is strong, and gravity,
$-\nabla\Phi$, is the gradient of that delay.

\paragraph{Light deflection.}
A light ray is itself a wave in $\Psi$.  Passing a source mass, it
propagates through the source's wave and interferes with it; the
cross-term makes its local phase velocity depend on position through
$\Phi$.  A transverse gradient of $\Phi$ pivots the wavefront, deflecting
the ray toward the mass --- the same gradient that accelerates a test mass
in free fall.  Light bends because a wave, like a clock, carries the
interference cross-term.  The refraction geometry is developed in
\S\ref{sec:refraction}.

\paragraph{Gravitational redshift.}
Redshift is time dilation read along a ray.  A wave emitted at radius
$r_1$, where the potential is $\Phi_1$, and received at $r_2$, where it is
$\Phi_2$, carries the emitter's tick rate to the receiver; the fractional
frequency change is
\begin{equation}
    \frac{\Delta\nu}{\nu} = \frac{\Phi_1 - \Phi_2}{c^2}.
\end{equation}
A wave climbing out of a well ($\Phi_1 < \Phi_2$) loses frequency.  No new
mechanism is invoked --- only \eqref{eq:time_dilation} applied at the two
ends of a path.

\begin{gapbox}[Open: extension beyond leading order]
Equations \eqref{eq:Phi_newton}--\eqref{eq:time_dilation} are exact at
$O(\Phi/c^2)$, the order at which the interference cross-term
\eqref{eq:cross_term} is the leading contribution.  Extension to higher
orders requires keeping the cubic and higher terms of the logarithmic
expansion \eqref{eq:log_expansion} --- treating the LogSE fully nonlinearly
rather than at quadratic order --- and confirming that the result
reproduces the parametrised post-Newtonian expansion of general relativity
(the PPN coefficients $\beta$, $\gamma$).  This is a clean test: SSV makes a
definite prediction at each order, and any departure from
$\beta = \gamma = 1$ would be a structural prediction distinct from GR.
The calculation is non-trivial but finite, and is carried out in
Paper~VII-b.
\end{gapbox}

\section{Path Bending and the Refraction Picture}
\label{sec:refraction}

\emph{[Section under development --- to be written.]}  This section will
treat the path-bending geometry in full: wavefront refraction in a graded
time delay, the bending of a trajectory toward greatest delay as a
quantitative statement, the lensing integral for light grazing a mass, and
the relation between the refraction picture and the potential gradient
$-\nabla\Phi$.

\section{Universal Coupling and the Absence of a Graviton}
\label{sec:universal}

\subsection{Universal Coupling: Why Everything Falls}

The interference mechanism is universal because every defect in the
medium --- regardless of internal structure, charge, or topology ---
radiates a wave, and any wave interferes with any other.  The amplitude of
that wave is set by the breather's mass-energy alone --- the fractional
volume excursion is $\delta V \propto m/\rho_0$ (Paper~II) --- so the
interference force scales as $m_1 m_2/r^2$ for any pair of objects.

This is the mechanical origin of the equivalence principle: the same
quantity (mass-energy) determines both the inertial response to
acceleration and the coupling to gravity itself, because both are
manifestations of the object's acoustic cross-section in the medium.  Every
other force is selective --- electromagnetism couples to chirality, the
strong force to topological charge, the weak force to core overlap --- but
gravity couples to the one property all defects share: they displace the
medium.

\subsection{Gravity Without a Graviton}

The SSV framework does not require a graviton.  The gravitational
potential $\Phi$ is the time-averaged interference of
\S\ref{sec:interference}; its $1/r$ long-range form is the spherical
dilution of a radiated wave, equivalently the Green's function of the
Laplacian.  What conventional quantum field theory calls graviton exchange
corresponds here to the interference of waves in the medium --- a collective
hydrodynamic process, not a particle interaction.

This is structurally consistent with the treatment of the other forces in
Paper~II: electromagnetism emerges from the transverse flow of chiral
vortices without postulating a separate photon field, and gravity emerges
from longitudinal density waves without postulating a separate graviton
field.  Both are modes of the single medium $\Psi$.

\paragraph{The ultraviolet problem dissolves.}
Conventional perturbative quantum gravity is non-renormalisable: when
gravity is treated as a quantum field theory in the standard way, the
loop integrals diverge at short distances and the infinities cannot
be removed by redefining a finite number of parameters.  The theory
has no built-in short-distance cutoff and breaks down at the Planck
scale, requiring replacement by an unknown theory of quantum gravity.

In SSV this problem does not arise --- not because we solve it, but
because the framework's domain of validity makes it impossible to
state.  SSV is a continuum effective theory valid down to the scale
of the heaviest stable defect,
\begin{equation}
    \ell_{\rm grain} \equiv a_p = \frac{\hbar}{m_p c}
    \approx 2.10\times10^{-16}\,\mathrm{m}.
    \label{eq:lgrain_cutoff}
\end{equation}
Below $\ell_{\rm grain}$, the LogSE continuum description ceases to
be self-consistent: no new stable structures exist, and the medium
supports only transient denting events (Paper~II).  Momentum
modes with $k \gg 1/\ell_{\rm grain}$ are outside the framework's
domain of validity --- not regulated away by hand, but simply absent
from the theory, in the same sense that Navier--Stokes has no
sub-molecular eddies.  The divergent integrals of perturbative quantum
gravity never appear because the modes that would generate them are
outside the scope of the description.

This is worth pausing on.  The ultraviolet catastrophe of quantum
gravity is usually treated as one of the deepest unsolved problems in
physics, motivating decades of work on string theory, loop quantum
gravity, and related programmes.  In SSV it is not solved --- it
simply does not arise.  The cutoff $\ell_{\rm grain}$ is not a new
postulate introduced to tame divergences; it is the healing length of
the heaviest stable defect, fixed by the proton mass alone.  The
same physical reasoning that gives Navier--Stokes its molecular
cutoff gives SSV its cutoff at $a_p$: at scales finer than the
smallest stable structure, the continuum description is silent.

Note the distinction from the electron healing length $\xi =
\hbar/(m_e c) \approx 3.86\times 10^{-13}\,\mathrm{m}$, the natural
scale for electron-defect physics but not the medium's ultraviolet
boundary: the true domain-of-validity boundary is $\ell_{\rm grain}
\approx \xi/1836$, set by the proton.

The absence of a graviton and the absence of UV divergences are the
same fact stated two ways: SSV has no gravitational field to quantise,
and therefore no quantisation problem to solve.

\section{From Update-Capacity Gradients to Emergent Geometry}
\label{sec:geometry}

The time-delay gradient of the preceding sections is the weak-field face of
a deeper statement: that geometry itself is a macroscopic limit of the
medium.  This section sketches the bridge; its full development --- the
effective metric to all orders, the post-Newtonian coefficients, and why
the standard quantum-gravity problem is mis-posed in SSV --- is the subject
of Paper~VII-b.

The curved-spacetime formulation of gravity emerges from the
Madelung equations in the hydrodynamic limit~\cite{Volovik,Barcelo}.
The acoustic metric
\begin{equation}
    g_{\mu\nu}^{\rm eff}
      = \frac{\rho_0}{m_0 c}
        \begin{pmatrix}
          -(c^2 - v^2) & -v_j \\
          -v_i         & \delta_{ij}
        \end{pmatrix}
    \label{eq:acoustic_metric}
\end{equation}
reduces to Minkowski when the background flow vanishes.  A single
interference potential $\Phi$ governs the metric: it sets both the
time-dilation factor $1+\Phi/c^2$ of \eqref{eq:time_dilation} and the
deflection of wave propagation.  That one potential does both --- the force
$-\nabla\Phi$ and the clock rate $1+\Phi/c^2$ --- is precisely the
weak-field structure of general relativity.  Gravitational lensing follows
automatically: light is a wave in the same $\Psi$, deflected by the same
interference gradient that bends the path of a test mass.  No separate
gravitational field is postulated at any point.

\section{Gravity in a Medium That Remembers}
\label{sec:remembers}

The treatment above is the mean-field story.  Two specific places in
the derivation depend on operations that the irreversibility framework of
Paper~III will revisit at finer resolution; they are flagged here so that
the present scope is honest.

\paragraph{Time-averaging the interference.}
Newton's constant was derived in Paper~II from the time-averaged
interference cross-term $\langle F\rangle = -\rho_0 \langle \dot V_1 \dot
V_2\rangle/(4\pi r^2)$.  The averaging brackets perform a classical
operation: integrate over many cycles of the proton breather oscillation at
$\omega_p = m_p c^2/\hbar$ to extract the static component.  In a
featureless medium this is unproblematic.  In a medium that records, what
gets averaged is itself history-dependent: each cycle deposits a small
amount of acoustic energy into the medium, and the breather of cycle $n+1$
propagates through a slightly different background than the breather of
cycle $n$.  The leading-order static result survives this refinement ---
the recoil of one cycle on the next is suppressed by the proton's sub-grain
coupling factor $(m_e/m_p)^3$, the same factor that makes $G$ small in the
first place --- but the corrections are non-zero and carry information about
the source's pulsation history.  These corrections are the
gravitational-wave memory effect (Christodoulou) read in this framework:
not as a peculiar property of asymptotic spacetimes, but as the direct
consequence of a medium that literally remembers.

\paragraph{The interference pattern as a long-time average.}
The static potential $\Phi$ is the long-time-averaged envelope of the
interfering pulsation waves.  In the framework's deeper picture,
$\rho(\vec x, t)$ is not a smooth deterministic function but a
coarse-grained description of an evolving record; the smoothness emerges
because reconnection events and their acoustic wakes are dense enough to
wash out at scales relevant to the equivalence-principle tests we currently
perform.  At sufficiently high precision and sufficiently short timescales,
this smoothing must fail: the gravitational potential at a point will
fluctuate around its mean value at a rate set by the local rate of
reconnection events.  No calculation of this fluctuation spectrum is given
here.  It is identified as the natural Paper~III follow-up to the present
mean-field gravity treatment.

\paragraph{What this does not change.}
\begin{sloppypar}
The Newton's-constant consistency check of Paper~II is unaffected: it
relates time-averaged quantities, all defined at the same level of coarse
graining.  The graviton-free argument is unaffected: there is still no
gravitational field to quantise.  The argument that $\ell_{\rm grain}$
supplies the natural ultraviolet cutoff is unaffected: it remains the
medium's minimum length.  What the forward connection makes explicit is only
that the time-averaging step inherits the medium's accumulated history at a
level the present treatment does not resolve --- and that this is not an
obstacle to the mean-field claims, only a scoping of where Paper~III takes
over.
\end{sloppypar}

\section{Discussion and Open Problems}
\label{sec:discussion}

The mechanism set out here is complete at leading order and introduces no
parameter beyond those fixed in Paper~I.  Gravity is the interference of the
pressure waves that breathers radiate into the one medium; the interference
slows local time; its gradient bends motion.  From the single potential
$\Phi$ the four classical gravitational phenomena follow without further
assumption, and the same logarithmic nonlinearity that makes the medium a
medium also makes gravity attractive and inverse-square.

What is established: the interference cross-term \eqref{eq:cross_term} and
its three structural properties; the four corollaries of
\S\ref{sec:consequences} at $O(\Phi/c^2)$; the absence of a graviton and of
an ultraviolet divergence.  What is deferred: the first-principles value of
the gravitational coupling --- the sub-grain factor that fixes $G$ ---
remains the open computation of Paper~II; the extension beyond leading
order, and with it the post-Newtonian coefficients and the full emergent
metric, is Paper~VII-b.  Two sections of the present paper are themselves
still open: \S\ref{sec:capacity}, the formal account of update capacity, and
\S\ref{sec:refraction}, the quantitative refraction geometry of path
bending.  Both are flagged in place and are the next drafting tasks.

\begin{thebibliography}{9}

\bibitem{Bjerknes1906}
V.~F.~K.~Bjerknes, \textit{Fields of Force}, Columbia University Press, 1906.

\bibitem{Volovik}
G.~E.~Volovik, \textit{The Universe in a Helium Droplet}, Oxford University
Press, 2003.

\bibitem{Barcelo}
C.~Barcel\'{o}, S.~Liberati, and M.~Visser, ``Analogue gravity,''
\textit{Living Rev.~Relativity} \textbf{14}, 3 (2011).

\end{thebibliography}

\end{document}
